\section{De los modelos de lenguaje al Agent System}

En el minuto 19, la clase pasa a los Agent. Este cambio no presenta otra arquitectura de modelo: coloca el reasoning procedure de la primera mitad dentro de un sistema que funciona de manera continua. El modelo debe recibir un objetivo, descomponer la tarea, invocar herramientas, observar resultados, guardar el estado y decidir si continúa. Al leer la segunda mitad hay que distinguir siempre entre \term{foundation model} y \term{agent system}: el primero aporta la capacidad de lenguaje y de razonamiento; solo el segundo se encarga del ciclo cerrado de acción.

\subsection{Agenda de Agent: del «por qué» al «cómo controlar»}

La diapositiva inicial enumera agent architecture, actions, computer interaction, long-term memory, personalization, multi-agent communication y future directions. El orden didáctico que propone es importante: primero explica por qué una sola llamada al modelo no basta, después aborda las interfaces de acción y solo al final llega a la colaboración en paralelo y a la gobernanza de la seguridad.

\lecturefigure{slide-28-agent-introduction.jpg}{Agenda de la mitad de la clase dedicada a Agent}{Video de clase de Stanford 00:19:51}

\begin{knowledgebox}{Lectura de la figura: los siete temas forman juntos el ciclo de vida del sistema}
Architecture determina los límites entre componentes; action interface determina qué puede hacerse; memory y personalization determinan el estado entre sesiones; multi-agent determina cómo ampliar el rendimiento y la especialización, y future directions trata reliability, permissions y sandboxing. Si solo se construye uno de estos módulos, no se obtiene un agent al que se le pueda delegar trabajo a largo plazo.
\end{knowledgebox}

\teachervoice{Div pregunta primero: «¿por qué no simplemente entrenar un modelo de lenguaje más grande?». Su respuesta no es que la capacidad del modelo carezca de importancia, sino que una llamada única no tiene chaining, recursion, persistent state ni environment feedback; las tareas reales requieren un sistema, no solo un monolith.}

\subsection{Por qué se necesitan los Agent: el lenguaje natural es solo la entrada}

La diapositiva Building AI Agents plantea una thesis de la clase de 2023: las personas expresan su intención en lenguaje natural y la IA opera el software y las máquinas. Esta visión reduce la carga de que el usuario haga clic uno por uno en la interfaz, pero «100 veces más rápido» y «ocurrirá en cinco años» son predicciones del ponente, no hechos que el curso haya verificado. Desde la ingeniería, una formulación más sólida es esta: una natural-language interface puede reducir el costo de expresar una tarea, y que aumente o no la eficiencia depende de la tasa de éxito, del número de confirmaciones y del costo de recuperarse de los fallos.

\lecturefigure{slide-29-building-ai-agents.jpg}{Building AI Agents: la visión del lenguaje natural a la operación de máquinas}{Video de clase de Stanford 00:20:30}

\begin{importantbox}{Definición de Agent como ciclo cerrado}
Sea el objetivo $g$, el estado interno $s_t$ y la observación del entorno $o_t$; en el paso $t$, el agent ejecuta
\[
p_t=\operatorname{Plan}(g,s_t,o_t),\qquad
a_t=\operatorname{Act}(p_t),\qquad
s_{t+1}=\operatorname{Update}(s_t,o_{t+1}).
\]
$p_t$ es el plan, $a_t$ es una acción sujeta a permisos y $o_{t+1}$ es la retroalimentación del entorno. Si el sistema solo genera $p_t$ sin ejecutar, observar ni actualizar, se parece más a un assistant o a un planner que a un agent loop completo.
\end{importantbox}

\begin{warningbox}{Automatizar la interfaz no elimina los costos por sí solo}
Un Agent puede reducir los clics manuales, pero introduce latencia de inferencia, costo en tokens, acciones erróneas, cambios en los sitios web, permisos de cuentas y necesidades de auditoría. Al evaluarlo, hay que comparar el end-to-end task success con la intervención humana total, no solo mostrar una grabación en la que todo sale bien.
\end{warningbox}

\subsection{Agent Stack: el modelo es solo uno de los componentes}

La diapositiva de arquitectura coloca el model en el centro y lo conecta con memory, planning, tools, reflection y environment. Aquí, Chain-of-Thought y Tree of Thoughts, vistos en la primera mitad, se convierten en métodos de planning; calculator, calendar y code interpreter se convierten en tools; el estado de corto y largo plazo pasa a memory, y la revisión posterior a la ejecución pasa a reflection.

\lecturefigure{slide-30-agent-stack.jpg}{Agent Stack con memory, planning, tools y reflection}{Video de clase de Stanford 00:22:24}

\begin{knowledgebox}{Términos clave: los seis componentes del Agent Stack}
\begin{tabularx}{\textwidth}{L{0.18\textwidth}>{\raggedright\arraybackslash}X >{\raggedright\arraybackslash}X}
\toprule
Componente & Problema que resuelve & Relación con esta clase \\
\midrule
Model & Entender el objetivo y generar planes candidatos & Es una neural compute unit; no acapara el estado del sistema \\
Planning & Convertir el objetivo en pasos con dependencias & Puede usar CoT, ToT, decomposition \\
Memory & Guardar el contexto actual y hechos entre tareas & Distingue context de persistent store \\
Tools & Ejecutar operaciones precisas o externas & calculator, browser, code interpreter, apps \\
Reflection & Revisar resultados y provocar revisiones & Requiere un verifier o señales observables del entorno \\
Environment & Devolver cambios reales de estado & Determina si el plan sigue siendo coherente con la realidad \\
\bottomrule
\end{tabularx}
\end{knowledgebox}

\teachervoice{La clase compara construir un agent con construir una computer: el LLM es como la CPU, pero también hacen falta RAM, almacenamiento a largo plazo, I/O, red, interfaz y configuración del usuario. El valor didáctico de esta analogía es impedir que se equipare directamente un benchmark del modelo con la capacidad del sistema.}

\subsection{MultiOn: un prototipo de clase, no una garantía general}

La diapositiva de MultiOn presenta el browser agent que el ponente construía en ese momento y usa «poder operar internet mediante lenguaje natural» para motivar las demos siguientes. Esta diapositiva debe leerse como metadatos sobre la fuente del prototipo de la clase de 2023: indica de qué sistema provienen las demos, no que todos los sitios web, cuentas y tareas tengan la misma confiabilidad. Esta sección establece primero los niveles de evidencia; el siguiente capítulo analiza cuadro por cuadro lo que el prototipo mostró realmente.

\lecturefigure{slide-31-multion-api.jpg}{Presentación del prototipo de browser-agent MultiOn}{Video de clase de Stanford 00:23:03}

\begin{importantbox}{Cuatro niveles de evidencia}
\begin{enumerate}
  \item \textbf{Imagen de la clase}: la grabación muestra efectivamente varios estados de una ejecución concreta del flujo;
  \item \textbf{Afirmación del ponente}: por ejemplo, «cero intervención humana» o «tiene permisos de compra», que son capability claims de la clase;
  \item \textbf{Verificación independiente}: requiere un conjunto de tareas reproducible, registros de fallos, versiones y un protocolo de evaluación;
  \item \textbf{Conclusión de despliegue}: además requiere evidencia de seguridad, permisos, recuperación y cumplimiento normativo.
\end{enumerate}
Estas notas atribuyen a la clase solo los dos primeros niveles y no los elevan automáticamente al tercero ni al cuarto.
\end{importantbox}

\subsection{Resumen de la sección}

Un Agent no puede resumirse con la definición «un modelo grande que invoca herramientas»: es un ciclo cerrado que forman en conjunto model, planning, memory, tools, reflection, environment y policy. Las live demos que siguen muestran con estados concretos el potencial de acción, y también dejan ver por qué las acciones irreversibles y las restricciones del entorno deben incorporarse al diseño.
